\documentclass[10pt]{article}
\usepackage[margin=0.76in]{geometry}
\usepackage{amsmath,amssymb}
\usepackage{microtype}
\usepackage{xcolor}
\usepackage{hyperref}
\hypersetup{colorlinks=true,urlcolor=blue!55!black}
\title{The planar Tingley problem is answered\\
a literature-status note for arXiv:2005.04794}
\author{Literature identification packet}
\date{13 August 2026}

\setlength{\parskip}{3pt}

\begin{document}
\maketitle

\begin{center}
\fbox{\parbox{0.9\linewidth}{\textbf{Status: literature already answered.}
The introductory two-dimensional open statement in arXiv:2005.04794 was
explicitly resolved by Taras Banakh, arXiv:2103.09268, Theorem 1.3.}}
\end{center}

\section*{Original statement}

Cueto-Avellaneda and Peralta state in the first paragraph of the introduction
to \emph{Can one identify two unital JB$^*$-algebras by the metric spaces
determined by their sets of unitaries?} (arXiv:2005.04794, page 1) that
Tingley's problem ``remains open even for two dimensional spaces.''

In precise terms, the planar problem asks: if $X$ and $Y$ are two-dimensional
Banach spaces and
\[
 f:S_X\longrightarrow S_Y
\]
is a surjective isometry between their unit spheres, must $f$ extend to a
surjective real-linear isometry $X\to Y$?  Equivalently, does every
two-dimensional Banach space have the Mazur--Ulam property?

\section*{Decisive later answer}

Banakh's \emph{Every 2-dimensional Banach space has the Mazur--Ulam property}
(arXiv:2103.09268) restates Tingley's question as Problem 1.1.  Its main result
is:
\begin{quote}
\textbf{Theorem 1.3.} Every 2-dimensional Banach space has the Mazur--Ulam
property.
\end{quote}
The abstract explicitly says that every isometry between the unit spheres of
two-dimensional Banach spaces extends to a linear isometry and that this
resolves Tingley's problem in dimension two.  Hence it is an exact affirmative
answer to the source's introductory open statement.

\section*{Provenance and scope}

This packet records a direct later-literature answer, not an independently
derived result. The supporting author explicitly identifies the theorem as a
resolution of the same problem. The identification is supported by the exact
problem formulation, theorem label, and matching dimensional scope.

The source paper's title question, concerning isometries of unitary sets of
unital JB$^*$-algebras, is answered inside arXiv:2005.04794 itself and is not
the signal recorded here. Banakh's theorem resolves only the two-dimensional
case; the unrestricted Tingley problem is not claimed to be solved by this
packet.

\section*{References}

M. Cueto-Avellaneda and A. M. Peralta,
\emph{Can one identify two unital JB$^*$-algebras by the metric spaces
determined by their sets of unitaries?}, arXiv:2005.04794.
\url{https://arxiv.org/abs/2005.04794}.

T. Banakh, \emph{Every 2-dimensional Banach space has the Mazur--Ulam
property}, arXiv:2103.09268, Theorem 1.3.
\url{https://arxiv.org/abs/2103.09268}.

\end{document}
